\documentclass[11pt,a4paper]{article}
\usepackage[margin=0.9in]{geometry}
\usepackage{amsmath,amssymb,booktabs,longtable,array,graphicx,float,placeins,siunitx}
\usepackage[numbers,sort&compress]{natbib}
\usepackage{xurl}
\usepackage[colorlinks=true,linkcolor=blue,citecolor=blue,urlcolor=blue]{hyperref}
\sisetup{detect-all}
\setlength{\emergencystretch}{2em}
\renewcommand{\thesection}{S\arabic{section}}
\renewcommand{\thesubsection}{S\arabic{section}.\arabic{subsection}}
\renewcommand{\thetable}{S\arabic{table}}
\renewcommand{\thefigure}{S\arabic{figure}}
\makeatletter
\renewcommand*\l@section{\@dottedtocline{1}{1.5em}{3.0em}}
\makeatother
\newcommand{\bestgraphic}[2][]{\IfFileExists{#2.pdf}{\includegraphics[#1]{#2.pdf}}{\includegraphics[#1]{#2.png}}}

\title{Supplementary Information\\Protocol Sensitivity in Remaining-Useful-Life Model Comparisons: Retrospective C-MAPSS Analysis and a Pre-Outcome-Frozen Battery Study}
\author{Zhihao Liu\\\small Independent Researcher, China}
\date{}

\begin{document}
% v5.5.0 release identity is recorded in the external submission sidecar.
\newcommand{\EvidenceReleaseID}{v5.5.0}
\newcommand{\EvidenceProtocolName}{ress-full-5x10-observed-prefix-analysis}
\newcommand{\EvidenceProtocolVersion}{3.2}
\newcommand{\EvidenceSourceReleaseTag}{v5.5.0}

\InputIfFileExists{generated/submission_identity.tex}{}{%
  \newcommand{\EvidenceTaggedCommit}{recorded in submission_identity.json}%
  \newcommand{\EvidenceTagTree}{recorded in submission_identity.json}%
  \newcommand{\EvidenceSourceContentCommit}{recorded in submission_identity.json}%
  \newcommand{\EvidenceRepositoryURL}{https://github.com/Ethan0infinity/rul-benchmark-protocol-audit}%
  \newcommand{\EvidenceTaggedCommitShort}{see submission identity record}%
  \newcommand{\EvidenceTagTreeShort}{see submission identity record}%
  \newcommand{\EvidenceSourceContentCommitShort}{see submission identity record}%
  \newcommand{\EvidenceManuscriptHash}{generated at release build}%
  \newcommand{\EvidenceSupplementHash}{generated at release build}%
}
\maketitle

\noindent\textbf{Evidence identity.} The source release is identified by the annotated GitHub tag \texttt{\EvidenceSourceReleaseTag}; its resolved tag, commit, and tree hashes are recorded in the external submission identity record. The v5.5.0 evidence archive is publicly available as a version-named file in the OSF project at \url{https://osf.io/download/6ab6a6735205d2cb18154961/}. Its SHA-256 is \nolinkurl{A390CA0216E9CA1238420CA3F76AAFDF536A18B07D1B0FF2631FAA7FDF50566B}. This is a public project file, not a new OSF registration or immutable version DOI. The final PDF hashes are listed in the same identity record; the OSF archive hash identifies the evidence ZIP, not the compiled PDFs.

\tableofcontents
\clearpage

\section{Scope, direction convention, and evidence vocabulary}

Every difference is first named configuration minus second named configuration. Negative values favor the first configuration for RMSE, NASA/engine, late-prediction ratios, and normalized maintenance-decision cost. ``Completed levels'' are fixed study objects, not draws from a defined population. ``Diagnostic'' means that an output describes the completed design without calibrated confirmatory coverage or a prospective decision role. Table~\ref{tab:evidence-vocabulary} defines the evidence vocabulary used in this concise Supplement.

\begin{table}[H]
\centering
\small
\caption{Evidence vocabulary used in both documents.}
\label{tab:evidence-vocabulary}
\begin{tabular}{@{}p{0.24\textwidth}p{0.31\textwidth}p{0.36\textwidth}@{}}
\toprule
Term & Operational meaning & Prohibited interpretation \\
\midrule
Fixed-task estimate & Paired mean and sign pattern over five completed coupled levels & Population effect, confirmatory test, or architecture main effect \\
Sensitivity & Observed direction, magnitude, or support over declared post-result choices & Prevalence over all defensible protocols \\
Diagnostic envelope & Resampling visualization of the completed design & Calibrated 95\% confidence interval \\
Prospectively frozen worked audit & Battery protocol timestamped before training and test-label evaluation & Independent-team replication or model superiority \\
Static RUL trigger threshold & Service/no-service mapping under declared RUL trigger threshold and cost grids & Optimized policy, monetary utility, or system reliability \\
\bottomrule
\end{tabular}
\end{table}

\section{Traceability scaffold and neighboring reporting frameworks}

The six-part traceability scaffold is supporting infrastructure for the protocol-sensitivity analysis, not a validated intervention or replacement for existing reporting frameworks. REFORMS is a consensus checklist for ML-based science \citep{Kapoor2024REFORMS}; Model Cards document model-level intended use and evaluation \citep{Mitchell2019ModelCards}; Eval Factsheets organize evaluation documentation \citep{Bordes2025EvalFactsheets}; and Evaluation Cards provide an interpretive evaluation-reporting layer with large-scale deployment evidence \citep{Ghosh2026EvaluationCards}. Chen et al.'s Evidence-Ledger Adjudication links claims to evidence packets and support relations for AI-assisted writing \citep{Chen2026EvidenceLedger}; its unit is a claim-verification packet, not an RUL benchmark comparison or protocol-sensitivity record. The present study has no independent comparative-effect result for the scaffold. Table~\ref{tab:construct-map} maps its fields to neighboring reporting frameworks without claiming superiority.

Figure~\ref{fig:supp-traceability} shows the scaffold, and Table~\ref{tab:supp-sixpart} defines its six operational fields.

\begin{table}[H]
\centering
\scriptsize
\setlength{\tabcolsep}{2.5pt}
\caption{Construct map. ``Explicit'' describes a named field, not evidence that the format is superior.}
\label{tab:construct-map}
\begin{tabular}{@{}p{0.17\textwidth}p{0.18\textwidth}p{0.17\textwidth}p{0.19\textwidth}p{0.20\textwidth}@{}}
\toprule
Framework & Primary object & Design timing & Provenance & Claim/wording constraint \\
\midrule
Six-part traceability scaffold & Individual benchmark claim & Explicit field & Explicit artifact field & Explicit permitted-wording field; efficacy untested \\
REFORMS & ML-based scientific study & Study-design recommendations & Computational-reproducibility module & Scope, metric, and limitation guidance \\
Evaluation Cards & Evaluation record and reader interpretation & Timing-relevant metadata & Central interpretive signal & Reader modes and comparability signals \\
Eval Factsheets & Evaluation documentation object & Context and procedure fields & Documentation-centered & Interpretive documentation across five dimensions \\
Evidence-Ledger Adjudication & Claim/evidence packet & Per-claim adjudication & Evidence packet and support relation & Routes unsupported, contradicted, or mixed claims for review \\
Model Cards & Trained model report & Training/evaluation procedure & Model/version record & Intended use, limitations, and disaggregated evaluation \\
Preregistration & Study plan & Central purpose & Immutable pre-outcome plan & Constrains post-result interpretation indirectly \\
Multiverse reporting & Analysis-choice distribution & Prospective or retrospective & Implementation-dependent & Exposes completed specification behavior \\
\bottomrule
\end{tabular}
\end{table}
\setlength{\tabcolsep}{6pt}

\begin{figure}[H]
\centering
\bestgraphic[width=0.90\linewidth]{figures/round19_governance_workflow}
\caption{Six-part traceability scaffold used as supporting documentation for the protocol-sensitivity analyses. The scaffold records comparison unit, design timing, evidence support, artifacts, completed sensitivities, and warranted wording; it does not upgrade the underlying evidence.}
\label{fig:supp-traceability}
\end{figure}

\begin{table}[H]
\centering
\small
\caption{Operational fields in the six-part traceability scaffold.}
\label{tab:supp-sixpart}
\begin{tabular}{@{}p{0.24\textwidth}p{0.67\textwidth}@{}}
\toprule
Field & Required content \\
\midrule
Comparison and unit & Quantity, direction, aggregation rule, and independent evaluation unit. \\
Design timing & When consequential choices were fixed and whether outcomes had already been inspected. \\
Evidence support & What the completed design can support without upgrading diagnostics to population inference. \\
Artifacts & Resolvable data rows, code, manifests, stable identifiers, and hashes needed to inspect the claim. \\
Completed sensitivities & Adverse, null, mixed, and reversing outcomes, reported together with their scale. \\
Claim wording & The strongest sentence warranted by the preceding five fields, without unsupported superiority or practical-benefit claims. \\
\bottomrule
\end{tabular}
\end{table}

\section{C-MAPSS finite-design evidence}

FD004 is development-informed because it supplied the original preference-development split. The statistical unit for endpoint comparisons is the matched official test engine; the five composite levels jointly change split, initialization, shuffle, and augmentation.

The archived raw files used for FD004 contain 249 unique training-unit identifiers and 248 test-unit identifiers. Their SHA-256 hashes are \nolinkurl{27ef6160b6a1dcb2613a88de9c239f763b223f02cdc41dc5cdedc5dc189b6218} (\texttt{train\_FD004.txt}) and \nolinkurl{1dc675fff0624bac10786927c6715b37d1297657137400d2b1a3138d777a3ba5} (\texttt{test\_FD004.txt}); both match the corresponding files inside the archived \texttt{CMAPSSData.zip} (SHA-256 \nolinkurl{74bef434a34db25c7bf72e668ea4cd52afe5f2cf8e44367c55a82bfd91a5a34f}). NASA's current Open Data landing page reports 248 training and 249 test trajectories \citep{NASA_CMAPSS_OpenData}. The experiments use the archived local files identified by these hashes.

\begin{table*}[!t]
\centering
\small
\setlength{\tabcolsep}{4.5pt}
\caption{Task-specific finite-design effects for OCM-Asym minus RAST-GRU. Ranges and signs summarize five completed coupled levels; they are not confidence intervals or population-inference statistics. RMSE effects are also shown as a percentage of the 125-cycle cap, and LPR effects as percentage-point changes; neither scale is an engineering equivalence margin. $N_{30}$ is the LPR denominator.}
\label{tab:primary-effects-only}
\begin{tabular}{llrrcrr}
\toprule
Task & Metric & Mean effect (descriptive scale) & Observed range & $-/+/0$ & $N_{test}$ & $N_{30}$ \\
\midrule
FD001 & RMSE (cycles) & -0.806 (-0.64\%) & [-2.453, +0.159] & 3/2/0 & 100 & 25 \\
FD001 & NASA / engine & -0.537 & [-1.305, +0.216] & 4/1/0 & 100 & 25 \\
FD001 & LPR@30 & -0.040 (-4.0 pp) & [-0.320, +0.160] & 3/2/0 & 100 & 25 \\
FD002 & RMSE (cycles) & -0.260 (-0.21\%) & [-0.891, +0.975] & 4/1/0 & 259 & 61 \\
FD002 & NASA / engine & -0.619 & [-1.137, +0.187] & 4/1/0 & 259 & 61 \\
FD002 & LPR@30 & -0.036 (-3.6 pp) & [-0.213, +0.082] & 2/3/0 & 259 & 61 \\
FD003 & RMSE (cycles) & +0.664 (+0.53\%) & [-0.422, +2.261] & 2/3/0 & 100 & 20 \\
FD003 & NASA / engine & +0.982 & [-0.133, +3.064] & 1/4/0 & 100 & 20 \\
FD003 & LPR@30 & +0.040 (+4.0 pp) & [-0.100, +0.200] & 2/2/1 & 100 & 20 \\
FD004 & RMSE (cycles) & +0.893 (+0.71\%) & [-1.143, +2.461] & 1/4/0 & 248 & 53 \\
FD004 & NASA / engine & -1.241 & [-9.647, +1.199] & 1/4/0 & 248 & 53 \\
FD004 & LPR@30 & -0.117 (-11.7 pp) & [-0.245, -0.038] & 5/0/0 & 248 & 53 \\
\bottomrule
\end{tabular}
\end{table*}

\begin{table}[H]
\centering
\small
\caption{Integer support for the targeted LPR@30 endpoint. Entries are observed late-event-count ranges across five coupled levels, shown as OCM-Asym / RAST-GRU.}
\label{tab:lpr30-event-support}
\begin{tabular}{lrrr}
\toprule
Task & Eligible $N_{30}$ & $\epsilon=0$ events & $\epsilon=5$ events \\
\midrule
FD001 & 25 & 0--4 / 0--8 & 0--0 / 0--0 \\
FD002 & 61 & 3--8 / 1--16 & 0--0 / 0--2 \\
FD003 & 20 & 1--6 / 1--5 & 0--1 / 0--1 \\
FD004 & 53 & 7--17 / 15--23 & 0--4 / 3--7 \\
\bottomrule
\end{tabular}
\end{table}


\section{Full split--stream crossing}

The 400-run extension pairs five split levels with ten training streams for each of four tasks and two complete configurations. Each task therefore has 50 matched split--stream cells and one estimate per cell. The formal execution began before the plan and Amendment 01 were frozen on 22 September 2026. The lock is pre-aggregation, not pre-execution preregistration. The complete aggregate was generated after the 400-row manifest was checked and hashed. All rows used validation last-risk-score checkpoint selection, an 80-epoch cap, and patience 12. The 400 runs stopped by the declared early-stopping rule between epochs 13 and 61; each history reached the patience condition and retained its selected checkpoint.

The heatmaps below show cell-level contrasts as OCM-Asym minus RAST-GRU. Negative cells favor OCM-Asym for the displayed error endpoint. Rows are split levels and columns are training streams. Each panel shows the finite design, not a sample from a split or stream population.

\begin{figure}[H]\centering
\begin{minipage}{0.48\linewidth}\centering\bestgraphic[width=\linewidth]{figures/full_5x10/heatmap_FD001_rmse}\\[-1mm]{\small FD001}\end{minipage}\hfill
\begin{minipage}{0.48\linewidth}\centering\bestgraphic[width=\linewidth]{figures/full_5x10/heatmap_FD002_rmse}\\[-1mm]{\small FD002}\end{minipage}\\[2mm]
\begin{minipage}{0.48\linewidth}\centering\bestgraphic[width=\linewidth]{figures/full_5x10/heatmap_FD003_rmse}\\[-1mm]{\small FD003}\end{minipage}\hfill
\begin{minipage}{0.48\linewidth}\centering\bestgraphic[width=\linewidth]{figures/full_5x10/heatmap_FD004_rmse}\\[-1mm]{\small FD004}\end{minipage}
\caption{RMSE contrasts across the 5-by-10 crossing.}
\label{fig:supp-crossing-rmse}
\end{figure}

\begin{figure}[H]\centering
\begin{minipage}{0.48\linewidth}\centering\bestgraphic[width=\linewidth]{figures/full_5x10/heatmap_FD001_nasa_per_engine}\\[-1mm]{\small FD001}\end{minipage}\hfill
\begin{minipage}{0.48\linewidth}\centering\bestgraphic[width=\linewidth]{figures/full_5x10/heatmap_FD002_nasa_per_engine}\\[-1mm]{\small FD002}\end{minipage}\\[2mm]
\begin{minipage}{0.48\linewidth}\centering\bestgraphic[width=\linewidth]{figures/full_5x10/heatmap_FD003_nasa_per_engine}\\[-1mm]{\small FD003}\end{minipage}\hfill
\begin{minipage}{0.48\linewidth}\centering\bestgraphic[width=\linewidth]{figures/full_5x10/heatmap_FD004_nasa_per_engine}\\[-1mm]{\small FD004}\end{minipage}
\caption{NASA/engine contrasts across the 5-by-10 crossing.}
\label{fig:supp-crossing-nasa}
\end{figure}

\begin{figure}[H]\centering
\begin{minipage}{0.48\linewidth}\centering\bestgraphic[width=\linewidth]{figures/full_5x10/heatmap_FD001_lpr30}\\[-1mm]{\small FD001}\end{minipage}\hfill
\begin{minipage}{0.48\linewidth}\centering\bestgraphic[width=\linewidth]{figures/full_5x10/heatmap_FD002_lpr30}\\[-1mm]{\small FD002}\end{minipage}\\[2mm]
\begin{minipage}{0.48\linewidth}\centering\bestgraphic[width=\linewidth]{figures/full_5x10/heatmap_FD003_lpr30}\\[-1mm]{\small FD003}\end{minipage}\hfill
\begin{minipage}{0.48\linewidth}\centering\bestgraphic[width=\linewidth]{figures/full_5x10/heatmap_FD004_lpr30}\\[-1mm]{\small FD004}\end{minipage}
\caption{LPR@30 contrasts across the 5-by-10 crossing.}
\label{fig:supp-crossing-lpr}
\end{figure}

\begin{table}[H]
\centering
\scriptsize
\setlength{\tabcolsep}{3.5pt}
\caption{Finite-design OCM-Asym minus RAST-GRU contrasts across the full split--stream crossing. Each task has 50 matched cells. Sign counts are negative/tied/positive.}
\label{tab:supp-full-crossing}
\begin{tabular}{@{}lrrrrrr@{}}
\toprule
Task & RMSE mean & RMSE $-/0/+$ & NASA/engine mean & NASA $-/0/+$ & LPR@30 mean & LPR $-/0/+$ \\
\midrule
FD001 & $-0.318$ & 33/0/17 & $-0.307$ & 36/0/14 & $-0.041$ & 28/8/14 \\
FD002 & $+0.171$ & 24/0/26 & $-0.124$ & 27/0/23 & $-0.031$ & 27/3/20 \\
FD003 & $+0.719$ & 18/0/32 & $+0.618$ & 22/0/28 & $-0.039$ & 22/13/15 \\
FD004 & $+0.653$ & 18/0/32 & $+0.203$ & 21/0/29 & $-0.113$ & 37/3/10 \\
\bottomrule
\end{tabular}
\end{table}

\begin{table}[H]
\centering
\small
\caption{Integer support for LPR@30 in the full split--stream crossing. Event ranges are the minimum--maximum count over the 50 cells per task.}
\label{tab:full-crossing-lpr-support}
\begin{tabular}{@{}lrrr@{}}
\toprule
Task & Matched cells & Eligible engines per cell & RAST/OCM late-event count range \\
\midrule
FD001 & 50 & 25 & 0--14 / 0--9 \\
FD002 & 50 & 61 & 0--23 / 0--11 \\
FD003 & 50 & 20 & 0--12 / 0--8 \\
FD004 & 50 & 53 & 5--30 / 3--23 \\
\bottomrule
\end{tabular}
\end{table}

\begin{longtable}{@{}llrrr@{}}
\caption{Pairwise finite-design rank support over the full split--stream crossing. Counts are OCM-Asym lower metric / tied / RAST-GRU lower metric among 50 matched cells.}\label{tab:supp-full-rank-support}\\
\toprule Task & Endpoint & OCM lower & Tie & RAST lower \\
\midrule
\endfirsthead
\toprule Task & Endpoint & OCM lower & Tie & RAST lower \\
\midrule
\endhead
FD001 & RMSE & 33 & 0 & 17 \\
FD001 & NASA/engine & 36 & 0 & 14 \\
FD001 & LPR@30 & 28 & 8 & 14 \\
FD001 & MAE & 26 & 0 & 24 \\
FD001 & SLPR@30,5 & 8 & 35 & 7 \\
FD001 & SLPR@30,10 & 0 & 50 & 0 \\
FD002 & RMSE & 24 & 0 & 26 \\
FD002 & NASA/engine & 27 & 0 & 23 \\
FD002 & LPR@30 & 27 & 3 & 20 \\
FD002 & MAE & 20 & 0 & 30 \\
FD002 & SLPR@30,5 & 6 & 35 & 9 \\
FD002 & SLPR@30,10 & 0 & 49 & 1 \\
FD003 & RMSE & 18 & 0 & 32 \\
FD003 & NASA/engine & 22 & 0 & 28 \\
FD003 & LPR@30 & 22 & 13 & 15 \\
FD003 & MAE & 13 & 0 & 37 \\
FD003 & SLPR@30,5 & 14 & 27 & 9 \\
FD003 & SLPR@30,10 & 2 & 48 & 0 \\
FD004 & RMSE & 18 & 0 & 32 \\
FD004 & NASA/engine & 21 & 0 & 29 \\
FD004 & LPR@30 & 37 & 3 & 10 \\
FD004 & MAE & 12 & 0 & 38 \\
FD004 & SLPR@30,5 & 41 & 3 & 6 \\
FD004 & SLPR@30,10 & 21 & 25 & 4 \\
\bottomrule
\end{longtable}


\section{Endpoint roles and development relations}

The main text distinguishes target-monitored endpoints from independent downstream validation signals. Table~\ref{tab:endpoint-role} records how the reported quantities relate to training and checkpoint selection. Reporting an endpoint on held-out test engines does not make it independent of model development or checkpoint selection.

\begin{table}[H]
\centering
\small
\caption{Endpoint-role map used to constrain interpretation.}
\label{tab:endpoint-role}
\begin{tabular}{@{}p{0.18\textwidth}p{0.25\textwidth}p{0.21\textwidth}p{0.23\textwidth}@{}}
\toprule
Endpoint & Training-objective relation & Checkpoint relation & Permitted interpretation \\
\midrule
RMSE & Aligned with point-error accuracy & Direct component of $S_{\mathrm{val}}$ & Target-monitored accuracy, not an independent endpoint \\
LPR@30, SLPR@30,10 & Thirty-cycle life weighting is shared; OCM adds a smooth near-failure positive-error term & Direct components of $S_{\mathrm{val}}$ & Target-conditioned near-failure-overestimation endpoints \\
NASA/engine & Not directly optimized & Not in $S_{\mathrm{val}}$ & Out-of-target asymmetric-error sensitivity \\
Alternative $\tau$ or positive $\epsilon$ & Related to, but not identical to, the 30-cycle target & Not in $S_{\mathrm{val}}$ & Threshold or tolerance sensitivity \\
ZIMLE, CMLE, CVaR, $A_\kappa$ & Objective-adjacent error summaries & Not in $S_{\mathrm{val}}$ & Descriptive magnitude/tail sensitivity, not maintenance utility \\
\bottomrule
\end{tabular}
\end{table}

\section{Magnitude and maintenance-decision sensitivities}

Candidate magnitude bands were introduced after outcome access and are not stakeholder-validated smallest effects of interest. The static action is $a_i(h)=\mathbb I(\hat y_i\le h)$ and the reference due-state is $z_i(h)=\mathbb I(y_i\le h)$. Normalized loss is $\rho z_i(h)[1-a_i(h)]+[1-z_i(h)]a_i(h)$. It excludes inspection duration, replacement lead time, downtime, failure probability, and physical consequence.

\begin{table}[!t]
\centering
\caption{Transparent candidate magnitude-band sensitivity.}
\label{tab:candidate-magnitude-band}
\small
\begin{tabular}{llrrrrr}
\toprule
Domain & Metric & Margin & Mean contrast & Within & Negative beyond & Positive beyond \\
\midrule
C-MAPSS & RMSE & 0.5 & +0.1227 & 5/20 & 7 & 8 \\
C-MAPSS & RMSE & 1 & +0.1227 & 14/20 & 3 & 3 \\
C-MAPSS & RMSE & 2 & +0.1227 & 17/20 & 1 & 2 \\
C-MAPSS & LPR30 & 0.01 & -0.0383 & 1/20 & 12 & 7 \\
C-MAPSS & LPR30 & 0.025 & -0.0383 & 2/20 & 12 & 6 \\
C-MAPSS & LPR30 & 0.05 & -0.0383 & 8/20 & 8 & 4 \\
NASA battery & RMSE & 0.5 & +2.5455 & 1/5 & 0 & 4 \\
NASA battery & RMSE & 1 & +2.5455 & 2/5 & 0 & 3 \\
NASA battery & RMSE & 2 & +2.5455 & 2/5 & 0 & 3 \\
NASA battery & LPR & 0.01 & -0.0777 & 2/5 & 2 & 1 \\
NASA battery & LPR & 0.025 & -0.0777 & 2/5 & 2 & 1 \\
NASA battery & LPR & 0.05 & -0.0777 & 2/5 & 2 & 1 \\
\bottomrule
\end{tabular}
\begin{minipage}{0.98\linewidth}\footnotesize Notes: Negative/positive are first-minus-second contrasts beyond the displayed symmetric margin. Margins were introduced after outcome access and are not application-validated smallest effects of interest. The table separates sign sensitivity from magnitude sensitivity; it is not a formal equivalence test.\end{minipage}
\end{table}

\begin{table}[!t]
\centering
\caption{Exploratory static maintenance-triage sensitivity for OCM-Asym minus RAST-GRU.}
\label{tab:maintenance-decision}
\scriptsize
\setlength{\tabcolsep}{3.5pt}
\begin{tabular}{rrrrrr}
\toprule
RUL trigger threshold $h$ & Ratio $\rho$ & $\Delta C$/engine & Cost $-/0/+$ & Action disagree. & $\Delta$ false-late \\
\midrule
10 & 2 & +0.0017 & 8/4/8 & 0.0228 & -0.0011 \\
10 & 5 & -0.0017 & 10/3/7 & 0.0228 & -0.0011 \\
10 & 10 & -0.0073 & 11/3/6 & 0.0228 & -0.0011 \\
20 & 2 & -0.0030 & 11/3/6 & 0.0178 & -0.0014 \\
20 & 5 & -0.0072 & 11/2/7 & 0.0178 & -0.0014 \\
20 & 10 & -0.0142 & 11/2/7 & 0.0178 & -0.0014 \\
30 & 2 & +0.0033 & 4/8/8 & 0.0124 & +0.0002 \\
30 & 5 & +0.0039 & 6/7/7 & 0.0124 & +0.0002 \\
30 & 10 & +0.0048 & 7/6/7 & 0.0124 & +0.0002 \\
\bottomrule
\end{tabular}
\begin{minipage}{0.98\linewidth}\footnotesize Notes: Service is triggered when predicted RUL is no greater than the declared RUL trigger threshold. Cost equals $C_L$ for a deferred action when true RUL is within the threshold and $C_E$ for premature service otherwise. Each summary gives an equal-weight mean over four tasks and five composite levels. Negative differences favor OCM-Asym. This is a post-result static decision sensitivity, not an optimized fleet policy, failure-probability model, or field utility estimate.\end{minipage}
\end{table}


\subsection{Observed-prefix policy grid}

The policy grid uses only current and prior predictions to trigger actions. True RUL enters only cost scoring. Thresholds are 10, 20, 30, and 40 cycles; lead times are 0, 5, and 10 cycles; persistence requirements are 1, 2, and 3 windows; corrective-cost ratios are 5, 10, and 20. The fixed-age comparator triggers at cycle 100, and an always-corrective comparator is also retained. A trigger fixed before the observation cutoff can be scored after the cutoff. The analysis never infers a new post-cutoff trigger.

Unresolved engines remain in the denominator for resolution and unresolved rates. Conditional costs are means among resolved engines only. Paired OCM--RAST cost uses engines resolved under both configurations. Reference-only, comparison-only, both-resolved, and neither-resolved counts are kept for each grid cell. Table~\ref{tab:dynamic-policy-paired} summarizes paired setting-level cost differences; its setting counts are not independent observations and supports vary across settings.

\begin{table}[H]
\centering
\scriptsize
\setlength{\tabcolsep}{4pt}
\caption{Observed-prefix paired cost contrasts over the declared dynamic policy grid. Each task has 5,400 dependent policy/split/stream settings. Cost delta is OCM-Asym minus RAST-GRU among jointly resolved engines; negative values favor OCM-Asym.}
\label{tab:dynamic-policy-paired}
\begin{tabular}{@{}lrrrrr@{}}
\toprule
Task & Settings $-/0/+$ & Median cost delta & Observed range & Paired support min/median/max \\
\midrule
FD001 & 2262/36/3102 & $+0.00519$ & $[-18.880,\ 5.397]$ & 1/22.5/33 \\
FD002 & 1707/9/3684 & $+0.01243$ & $[-6.263,\ 3.493]$ & 18/60.5/92 \\
FD003 & 2079/27/3294 & $+0.00969$ & $[-18.885,\ 12.588]$ & 2/18/30 \\
FD004 & 1623/9/3768 & $+0.01406$ & $[-8.797,\ 4.183]$ & 7/51.5/83 \\
\bottomrule
\end{tabular}
\end{table}

\begin{table}[H]
\centering
\scriptsize
\setlength{\tabcolsep}{4pt}
\caption{Ranges across the 5,400 declared dynamic-policy settings per task for resolution rate and conditional mean cost. Costs are calculated among resolved engines; ranges are not confidence intervals.}
\label{tab:dynamic-policy-resolution}
\begin{tabular}{@{}llrr@{}}
\toprule
Task & Configuration & Resolution-rate range & Mean cost among resolved range \\
\midrule
FD001 & RAST-GRU & 0.010--0.360 & 1.097--20.000 \\
       & OCM-Asym & 0.060--0.360 & 1.117--5.851 \\
FD002 & RAST-GRU & 0.070--0.371 & 1.134--7.439 \\
       & OCM-Asym & 0.104--0.390 & 1.145--4.927 \\
FD003 & RAST-GRU & 0.020--0.360 & 1.085--20.000 \\
       & OCM-Asym & 0.040--0.400 & 1.105--13.708 \\
FD004 & RAST-GRU & 0.028--0.363 & 1.274--10.811 \\
       & OCM-Asym & 0.036--0.359 & 1.267--10.270 \\
\bottomrule
\end{tabular}
\end{table}


\section{Retrospective protocol sensitivities retained in the archive}

The concise document reports the main interpretation-changing checks; the complete row-level registries and all secondary tables remain machine readable. None is confirmatory. Table~\ref{tab:sensitivity-dashboard} provides a compact dashboard of those completed sensitivity families.

\begin{table}[H]
\centering
\small
\caption{Compact sensitivity dashboard. Directions are OCM-Asym minus RAST-GRU.}
\label{tab:sensitivity-dashboard}
\begin{tabular}{@{}p{0.23\textwidth}p{0.25\textwidth}p{0.42\textwidth}@{}}
\toprule
Family & Completed support & Interpretation \\
\midrule
Fixed-split streams & Ten optimization streams per task at one split & LPR mean is lower on all four tasks, but split variation is absent \\
Selected split--stream grid & Negative LPR directions: 4/9, 5/9, 4/9, 7/9 on FD001--FD004 & Result-informed finite grid; no split or stream population inference \\
FD004 preprocessing & Window 50 reverses mean RMSE and LPR; 21 sensors reverses mean LPR & Representation, capacity, and optimization exposure change together \\
LOFO & Random missingness is adverse; clean correction reverses some block/drift summaries & Conditional on normalized-coordinate augmentation construction \\
N-CMAPSS proxy & Same three test units under thinning/window controls & Representation dependence, not external validation \\
\bottomrule
\end{tabular}
\end{table}

\subsection{N-CMAPSS representation diagnostic}

The N-CMAPSS DS02-006 analysis is deliberately narrower than an external validation study \citep{AriasChao2021NCMAPSS}. It keeps every 100th record, forms overlapping length-50 windows, and reuses the same three official held-out test units while selected development, sampling, and window choices are changed. Because the representation is heavily thinned and the test units are unchanged, the analysis is used only to show representation dependence. It is not treated as field run-to-failure evidence, a CruiseBench reproduction, or an independent benchmark validation.

Under the declared validation-only condition-count rule, $K=4$ has the lowest validation risk (0.189), but retrospective test behavior is not uniformly favorable under the reused three test units. The result is therefore retained as a computational representation diagnostic rather than an external RUL validation claim.

\subsection{P10 comparator provenance}

The archived P10 extension is not part of the primary comparison. It contains 400 C-MAPSS runs, paired as 200 task--split--stream cells, together with grouped battery and baseline runs. Its C-MAPSS labels are \texttt{rast\_gru\_v2} and \texttt{rast\_gru}; the source record does not establish that \texttt{rast\_gru\_v2} is OCM-Asym. In addition, P10 selects checkpoints with \texttt{val\_last\_rmse}, whereas the focal C-MAPSS protocol uses \texttt{val\_last\_risk\_score}. The two result families therefore remain separate. P10 is retained here as provenance, not as a replication or pooled estimate.

\section{Latency correction and release identity}

The earlier manuscript/Supplement mismatch for the OCM-Asym CPU batch-1 value is closed by generating Table~\ref{tab:inference-latency-protocol} directly from \path{paper_outputs/release_v1/inference_latency_protocol_summary.csv}. The authoritative median is 1.612 ms, not 1.167 ms. The former Supplementary Table S27 is superseded by the following generated exhibit; no manually transcribed timing value is authoritative.

\begin{table}[!t]
\centering
\scriptsize
\caption{Repeated FD004 forward-pass timing. GPU measurements use 50 warm-up and 200 synchronized timed iterations; CPU measurements use 20 warm-up and 100 timed iterations. Preprocessing and host--device transfer are excluded. Values summarize five independently trained checkpoints.}
\label{tab:inference-latency-protocol}
\begin{tabular}{lrlrrrr}
\toprule
Point & Params & Device & Batch & Median ms & Checkpoint range & Median [p10,p90] \\
\midrule
OCM-Asym & 54,155 & CPU & 1 & 1.612 & [1.582, 1.999] & [1.389, 2.057] \\
OCM-Asym & 54,155 & GPU & 1 & 1.886 & [1.584, 2.317] & [1.484, 2.295] \\
OCM-Asym & 54,155 & GPU & 128 & 1.990 & [1.720, 2.159] & [1.603, 2.275] \\
RAST-GRU & 87,568 & CPU & 1 & 1.658 & [1.610, 1.723] & [1.547, 1.839] \\
RAST-GRU & 87,568 & GPU & 1 & 1.836 & [1.656, 2.076] & [1.511, 2.080] \\
RAST-GRU & 87,568 & GPU & 128 & 1.868 & [1.574, 2.665] & [1.440, 2.484] \\
\bottomrule
\end{tabular}
\end{table}


\section{Analysis chronology}

The timeline separates original training, result-informed reanalysis, external battery freezing, later artifact correction, and the present journal revision. Dates and clock times added for the full crossing are local Asia/Shanghai time unless a UTC suffix is printed. The 22 September entry records that execution was already underway when the plan was frozen; it does not assert pre-execution registration. A later public release can repair accessibility but cannot change a retrospective analysis into prospective evidence.

\small
\begin{longtable}{@{}>{\raggedright\arraybackslash}p{0.13\textwidth}>{\raggedright\arraybackslash}p{0.29\textwidth}>{\raggedright\arraybackslash}p{0.22\textwidth}>{\raggedright\arraybackslash}p{0.27\textwidth}@{}}
\caption{Complete analysis and release timeline. Later artifact repair does not change earlier design timing.}\label{tab:analysis-timeline} \\
\toprule
Date/period & Event & Source record & Evidence status \\
\midrule
\endfirsthead
\toprule
Date/period & Event & Source record & Evidence status \\
\midrule
\endhead
2026-07-04 & C-MAPSS project scaffold and initial experiment design & local project records & precedes later governance analyses; no external timestamp \\
before 2026-07-20 & Common-protocol C-MAPSS training grid completed & analysis chronology & results later inspected; retrospective evidence \\
2026-07-26 & Crossed seed-engine reanalysis and threshold diagnostics & analysis chronology & result-informed reanalysis; no retraining \\
2026-07-26/27 & Normalization, perturbation, mask, and latency audits & analysis chronology & post-result sensitivity families \\
2026-07-29 to 2026-08-08 & Retrospective extension and release-governance revisions & round 10--19 response records & completed after earlier result access \\
2026-08-09 & Evidence snapshot 2026-08-09-r21 assembled & workflow state & local evidence snapshot; not a prospective protocol \\
2026-08-12\\09:10:33Z & Battery protocol externally registered & external timestamp record & precedes battery training and test-label evaluation \\
after 2026-08-12 registration & Registered battery pair trained and evaluated & battery results audit & prospectively frozen worked execution; same investigator \\
2026-08-16 & Source provenance correction verified in v5.3.1 & Git tag v5.3.1 and repository history & corrected source snapshot; does not alter evidence timing \\
2026-08-16 & RESS major-revision exploratory analyses generated & major-revision evidence generator & post-result maintenance, equivalence, baseline, and target-definition sensitivities \\
2026-08-16 & RESS submission object rebuilt & Git tag v5.4.0 and submission\_identity.json & release synchronization and journal-format revision; no evidence-status upgrade \\
2026-08-16 & v5.4.0-p1 local submission-layer patch assembled & submission patch record & abstract reduced to journal limit; manifest self-hash corrected; selected-grid table transcription repaired; stale unused macros removed; no scientific result changed \\
2026-08-17 & v5.4.0-p2 local submission-compliance patch assembled & submission patch record & corresponding-author contact completed; formal OSF dataset citation added; vector Figure 5 added; PDF-first artwork and source-package bibliography resolution verified; no scientific result changed \\
2026-08-17 & v5.4.0-p3 privacy-safe submission patch assembled & submission patch record & residential contact isolated to confidential Editorial Manager metadata; public manuscript/source scrubbed; declaration mapping and portable test contract repaired; no scientific result changed \\
2026-08-20 & v5.4.0-p4 final migration-alignment patch assembled & submission patch record and migrated GitHub repository & migrated GitHub tag resolved to b48c7e302cb0...; repository URL/identity synchronized; figure/table text citations and evidence-manifest chronology repaired; no scientific result changed \\
2026-08-20 & v5.4.0-p5 submission-documentation and compatibility patch assembled & submission patch record & disclosure wording and document metadata normalized; no scientific result changed \\
2026-09-12 & v5.4.0-p6-readability readability revision assembled & submission patch record & abstract, introduction, section openings, discussion, and conclusion restructured for broad reliability-engineering readability; scientific result changed: no \\
2026-09-12 & v5.4.0-p7-editorial-polish editorial-polish revision assembled & submission patch record & title, abstract, front-loaded terminology, section headings, captions, and submission-facing descriptions polished for editor readability; scientific result changed: no \\
2026-09-12 & v5.4.0-p8-reviewer-risk-reduction scientific-narrative revision assembled & submission patch record & scientific contribution order, research questions, related-work emphasis, results/discussion flow, and submission-facing descriptions restructured to foreground RUL evaluation robustness and bounded decision relevance; scientific result changed: no \\
2026-09-12 & v5.4.0-p9-reviewer-hardening reviewer-hardening revision assembled & submission patch record & title reframed as protocol sensitivity; scientific contributions reduced to three; C-MAPSS rationale clarified; endpoint-role, N-CMAPSS, and reproducibility detail moved to Supplement; scientific result changed: no \\
2026-09-22 & Full 5-by-10 execution was underway before analysis plan and Amendment 01 were frozen & local launcher, plan, Amendment 01, and run manifest & pre-aggregation analysis lock; not pre-execution preregistration \\
2026-09-25 18:40 & Full 400-cell run manifest complete; configuration and artifact QA passed & canonical manifest and per-run outputs & 400 completed cells; all runs stopped under the declared patience rule \\
2026-09-25 19:06 & Frozen finite-design aggregation and pairwise rank-support tables generated & aggregation provenance and output hashes & 400 runs and 200 matched split--stream cells in total; descriptive only \\
2026-09-25 19:46 & Observed-prefix dynamic policy grid completed & policy protocol, summary, paired comparison, and engine-level outputs & post-result exploratory grid; no field or policy-benefit claim \\
\bottomrule
\end{longtable}

\normalsize


\section{Pre-outcome-frozen battery study details}

Discharge observations are indexed by $j=0,\ldots,J_i$. The trained target is $J_i-j$; $J_i+1$ is the observed discharge count. Four held-out batteries are the independent units. Five training streams are repeated fitting variations, and eligible windows are metric denominators rather than independent replicates.

\begin{table}[H]
\centering
\small
\caption{Prospectively frozen NASA battery evaluation. Values are the mean and observed range over five registered training streams. Differences are Battery-GRU-Asym minus Battery-GRU-Core. Counts describe stream-level signs and are not inferential probabilities.}
\label{tab:battery-primary}
\begin{tabular}{@{}lccc@{}}
\toprule
Configuration/contrast & Macro RMSE (cycles) & LPR$_{10\%}$ & Direction support \\
\midrule
Battery-GRU-Core & 30.051 [21.299, 33.687] & 0.8610 [0.5172, 1.0000] & 5 streams \\
Battery-GRU-Asym & 32.596 [22.166, 36.776] & 0.7833 [0.5909, 1.0000] & 5 streams \\
\midrule
Asym $-$ Core & +2.546 & -0.0777 & RMSE 0/0/5; LPR 2/2/1 \\
\bottomrule
\end{tabular}
\begin{minipage}{0.96\linewidth}\footnotesize Direction support is negative/zero/positive. The four test batteries contribute equally within every stream. The registered rule classifies the joint result as an observed accuracy--late-overprediction trade-off.\end{minipage}
\end{table}

\begin{table}[H]
\centering
\small
\caption{Battery-level support averaged over five registered streams. Each battery is an independent evaluation unit; window counts are denominators, not replicates.}
\label{tab:supp-battery-units}
\begin{tabular}{@{}lrrrrrrr@{}}
\toprule
Battery & Eligible & Core RMSE & Asym RMSE & $\Delta$RMSE & Core LPR & Asym LPR & $\Delta$LPR \\
\midrule
B0007 & 17 & 58.086 & 62.041 & +3.955 & 0.8471 & 1.0000 & +0.1529 \\
B0044 & 12 & 27.490 & 30.696 & +3.206 & 0.8333 & 0.5333 & -0.3000 \\
B0048 & 8 & 11.517 & 11.818 & +0.300 & 1.0000 & 1.0000 & +0.0000 \\
B0055 & 11 & 23.111 & 25.831 & +2.720 & 0.7636 & 0.6000 & -0.1636 \\
\bottomrule
\end{tabular}
\end{table}

\begin{table}[H]
\centering
\caption{Post-registration exploratory battery baselines under global training-only normalization.}
\label{tab:battery-baselines}
\scriptsize
\setlength{\tabcolsep}{4pt}
\begin{tabular}{p{0.34\linewidth}rrrr}
\toprule
Estimator & Macro RMSE & Macro LPR$_{10\%}$ & Pooled late/eligible & Pooled frequency \\
\midrule
training-life-prior & 23.720 & 0.2500 & 8.0/48.0 & 0.1667 \\
capacity-index-linear & 25.998 & 0.8824 & 40.0/48.0 & 0.8333 \\
window-summary-ridge & 12.711 & 0.9773 & 47.0/48.0 & 0.9792 \\
Battery-GRU-Core (five-stream mean) & 30.051 & 0.8610 & 40.8/48.0 & 0.8500 \\
Battery-GRU-Asym (five-stream mean) & 32.596 & 0.7833 & 38.0/48.0 & 0.7917 \\
\bottomrule
\end{tabular}
\begin{minipage}{0.98\linewidth}\footnotesize Notes: The three simple estimators were added after registration and do not enter the registered hypothesis or wording rule. GRU rows are five-training-stream means. LPR is an equal-battery macro average; pooled counts use 48 eligible windows per stream for the registered models and must not be reverse-calculated from macro LPR.\end{minipage}
\end{table}

\begin{table}[H]
\centering
\caption{Battery target-denominator definition sensitivity at the primary $q=0.10$, $\epsilon=0$ setting.}
\label{tab:battery-index-sensitivity}
\small
\begin{tabular}{lrrrr}
\toprule
Denominator & $\Delta$ macro RMSE & $\Delta$ macro LPR & Core/Asym pooled events & Eligible windows \\
\midrule
$J_i$ & +2.5455 & -0.0777 & 204/190 & 240 \\
$J_i+1$ & +2.5455 & -0.0777 & 204/190 & 240 \\
\bottomrule
\end{tabular}
\begin{minipage}{0.98\linewidth}\footnotesize Notes: $J_i$ is the final zero-based discharge index and $J_i+1$ is the observed discharge-count denominator. The trained target and predictions remain $J_i-j$; only the life-fraction denominator changes. RMSE is therefore invariant. Event counts are pooled over four batteries and five streams, whereas LPR is first averaged equally across batteries and then across streams.\end{minipage}
\end{table}


The registered global-normalization LPR is an equal-battery macro average. Pooled event counts over 48 eligible windows need not equal macro LPR multiplied by 48. Over the complete registered $(q,\epsilon)$ grid and both normalization rules, changing the life-fraction denominator from $J_i$ to $J_i+1$ shifts the Asym-minus-Core macro-LPR contrast by at most 0.0625 in absolute value; the primary cell is unchanged.

\section{Independent comparator-study status}

An executable study-ready package is located at \path{studies/six_link_comparator_study/}. It specifies a randomized blinded-label comparison, a target of at least 12 independent raters, balanced packet assignment, a frozen defect key, and rater-level outcomes for defect detection, false positives, claim-strength agreement, completion time, and burden. At this release the response file contains zero participants. The author is not counted as an independent rater, generative-AI reviewers do not satisfy that criterion, and recruitment is blocked pending a human-participant/ethics determination. The only permitted current conclusion is that two author-run worked records and a future comparator protocol exist; comparative efficacy is untested. Table~\ref{tab:comparator-status} records the readiness state and resulting claim boundary.

\begin{table}[H]
\centering
\small
\caption{Comparator-study readiness and claim boundary.}
\label{tab:comparator-status}
\begin{tabular}{@{}p{0.29\textwidth}p{0.24\textwidth}p{0.37\textwidth}@{}}
\toprule
Object & Current state & Consequence for this article \\
\midrule
Protocol and construct map & Prepared after current result inspection & May define a future study; no prospective status here \\
Assignment and response schemas & Present and validated & Demonstrates executability of data capture only \\
Independent rater responses & 0 & No effect size, agreement, timing, or burden result may be reported \\
Ethics determination & Not yet obtained & No recruitment is authorized by this package \\
Current scaffold evidence & Two same-author worked cases & Report as traceability support, not an effective intervention \\
\bottomrule
\end{tabular}
\end{table}

\section{Reproducibility checks and artifact record}

Reproducibility checks are reported at three distinct levels. \emph{Stored-result verification} checks machine-readable outputs, row counts, manifests, and hashes without retraining. \emph{Exhibit regeneration} rebuilds reader-facing tables and figures from declared source rows. \emph{Full or sentinel retraining} reruns a declared training path and is held to a prediction-level criterion. Passing one level does not imply passing the others, and none changes the original timing status of an analysis.

For C-MAPSS, the release record reports a clean stored-result verification and one same-stack FD001 seed-42 sentinel refit that exactly reproduced all 100 RAST-GRU and OCM-Asym predictions, including selected epochs and event summaries. The compact submission supports manifest, hash, schema, and stored-output checks but does not contain the full checkpoint/prediction tree needed to rerun that sentinel refit. The complete post-study evidence archive is available at OSF \url{https://osf.io/u76ze/}, DOI \nolinkurl{10.17605/OSF.IO/U76ZE}; a public source and verified-output layer is available at \url{https://github.com/Ethan0infinity/rul-benchmark-protocol-audit}. These records document bounded reproduction paths, not investigator-independent reproducibility or prospective C-MAPSS evidence.

For the battery study, the externally timestamped protocol ZIP and the locally executed protocol ZIP are byte-identical, with SHA-256 \nolinkurl{f192bb6bdd6188742bb7da4bf17d3fec164e5e2b5ccd87eb272dbb4f1dbd11fd}. Prediction, label, result, history, environment-deviation, and execution-hash records are retained separately after execution. The post-execution evidence object is included unchanged in the public OSF archive. These checks support two bounded statements: the registered protocol preceded outcome evaluation, and the completed evidence is inspectable. They do not provide investigator independence or show that the traceability scaffold improves reproducibility in other studies.

Selected claim-to-artifact navigation appears below. Complete paths, hashes, commands, and generator identities remain in the release records and OSF evidence archive.

\section{Artifact navigation and identity contract}

Table~\ref{tab:artifact-map} lists the core machine-readable sources used by this concise submission.

\begin{table}[H]
\centering
\scriptsize
\caption{Core machine-readable sources for the concise submission.}
\label{tab:artifact-map}
\begin{tabular}{@{}p{0.28\textwidth}p{0.30\textwidth}p{0.32\textwidth}@{}}
\toprule
Claim family & Source & Generator or verifier \\
\midrule
Fixed C-MAPSS effects & \path{paper_outputs/release_v1/primary_seed_level_effects.csv} & release-v1 exhibit generators \\
Full split--stream crossing & \path{paper_outputs/full_5x10_crossing/analysis/full_5x10_paired_contrasts.csv} & \path{scripts/aggregate_full_5x10_crossing.py} \\
Observed-prefix policy & \path{paper_outputs/full_5x10_crossing/dynamic_policy/dynamic_policy_paired_model_comparison.csv} & \path{scripts/run_dynamic_maintenance_policy.py} \\
Maintenance decisions & \path{paper_outputs/ress_major_revision/maintenance_decision_cell_level.csv} & stored-output verification; complete inputs required for regeneration \\
Candidate magnitude bands & \path{paper_outputs/ress_major_revision/*practical_equivalence.csv} & stored-output verification; complete inputs required for regeneration \\
Battery baselines and $J_i$ audit & \path{paper_outputs/ress_major_revision/battery_*.csv} & hashes verified locally; complete inputs required for regeneration \\
Comparator study & \path{studies/six_link_comparator_study/} & \path{scripts/prepare_comparator_study.py --check-only} \\
Submission identity & \path{submission_identity.json} & release build script and independent SHA-256 check \\
\bottomrule
\end{tabular}
\end{table}

The release identity uses one upstream annotated tag plus an external local-patch build record. The compact major-revision manifest inventories stored outputs without self-hashing; its generator requires the complete prediction and prepared-data inputs listed in that manifest for regeneration. A Git commit cannot contain its own cryptographic identifier without a self-reference paradox; therefore the tagged source contains the build logic and release tag, while the regenerated submission object records the resolved tag commit, tagged tree, source-content commit, final PDF hashes, archive hash, and timestamp in \texttt{submission\_identity.json}. This contract avoids presenting a prior commit, stale PDF, or mutable branch name as the identity of the submitted object.

The 5-by-10 crossing and observed-prefix policy outputs summarized here are included in the version-named OSF evidence archive listed above. The OSF project DOI identifies the project record; no separate version DOI is assigned to this file. GitHub source and compiled-document identities are recorded separately so that the evidence ZIP, source tree, and PDFs are not conflated.

\bibliographystyle{elsarticle-num-names}
\bibliography{references}
\end{document}
